% GENERATED FILE. Edit canonical lessons, then run npm run generate:books.
% canonical-source-hash: fnv1a64:01b1720c1611172b
% canonical-lessons: FA-C528-ketabforushi, FA-C528-onvan, FA-C528-jeld, FA-C528-nasher, FA-C528-motarjem

\chapter{Books}
\label{ch:fa-books}
\hlchaptermodality{528}

\begin{chapteropening}
\textbf{By the end of this chapter:} \emph{I can talk about books and the people who make them.}

Complete the last lesson of chapter 528: I can talk about books and the people who make them.
\end{chapteropening}

\section[\texorpdfstring{\fa{کتابفروشی} (ketâbforushi)}{کتابفروشی (ketâbforushi)}]{\fa{کتابفروشی} (ketâbforushi) — a bookshop}

\label{lesson:FA-C528-ketabforushi}



\begin{quote}
Before the new one: say the Persian for a dancer, then the Persian for to admire.
\end{quote}

\subsection*{You'll want to know: \fa{کتابفروشی}}
\textbf{\fa{کتابفروشی}} (\emph{ketâbforushi}) — \textquotedblleft{}a bookshop\textquotedblright{}.

\begin{quote}
\fa{نزدیک} \fa{خانه} \fa{ما} \fa{یک} \fa{کتابفروشی} \fa{قدیمی} \fa{هست}.
\end{quote}

\begin{quote}
\emph{nazdik-e khâne-ye mâ yek ketâbforushi-ye qadimi hast.}
\end{quote}

\begin{quote}
There's an old bookshop near our house.
\end{quote}

\emph{ketâb} 'book' + \emph{forush} 'selling' (stem of \emph{forukhtan}) + \emph{-i}.

\begin{cousinweb}
\emph{ketâbforushi} is built in Persian on \emph{ketâb}, 'book', from Arabic \emph{kitāb} (EtymDB), and \emph{forush}, the stem of native \emph{forukhtan}, 'to sell', with \emph{-i} for the shop.
\end{cousinweb}

\begin{grammarlens}[title={What you've built}]
One more word for telling, explaining and describing at B1.
\end{grammarlens}

\subsection*{Your turn}
\begin{itemize}
\raggedright
  \item \emph{Say it:} \emph{ketâbforushi}
  \item \emph{Say it:} \emph{ketâbforushi}, once more
  \item \emph{Say it:} a sentence of your own with \emph{ketâbforushi}
  \item \emph{Recall:} say the Persian for a dancer, then the Persian for to admire, then say \emph{ketâbforushi} again
\end{itemize}

\begin{tcolorbox}[breakable,colback=teal!4,colframe=teal!35!black,title={Before you move on}]
What does \fa{کتابفروشی} mean? (\textquotedblleft{}A bookshop\textquotedblright{}.) Say it once more.
\end{tcolorbox}
\section[\texorpdfstring{\fa{عنوان} ('onvân)}{عنوان ('onvân)}]{\fa{عنوان} ('onvân) — a title; a heading}

\label{lesson:FA-C528-onvan}



\begin{quote}
Before the new one: say the Persian for to admire, then the Persian for a bookshop.
\end{quote}

\subsection*{You'll want to know: \fa{عنوان}}
\textbf{\fa{عنوان}} (\emph{'onvân}) — \textquotedblleft{}a title; a heading\textquotedblright{}.

\begin{quote}
\fa{عنوان} \fa{این} \fa{کتاب} \fa{خیلی} \fa{جالب} \fa{است}.
\end{quote}

\begin{quote}
\emph{'onvân-e in ketâb kheyli jâleb ast.}
\end{quote}

\begin{quote}
The title of this book is very interesting.
\end{quote}

Plural \emph{'anâvin} (Arabic). \emph{be 'onvân-e} — as, in the role of.

\begin{cousinweb}
\emph{'onvân} was borrowed from Arabic \emph{ʿunwān} (EtymDB), a word for the title or address on a letter.
\end{cousinweb}

\begin{grammarlens}[title={What you've built}]
One more word for telling, explaining and describing at B1.
\end{grammarlens}

\subsection*{Your turn}
\begin{itemize}
\raggedright
  \item \emph{Say it:} \emph{'onvân}
  \item \emph{Say it:} \emph{'onvân}, once more
  \item \emph{Say it:} a sentence of your own with \emph{'onvân}
  \item \emph{Recall:} say the Persian for to admire, then the Persian for a bookshop, then say \emph{'onvân} again
\end{itemize}

\begin{tcolorbox}[breakable,colback=teal!4,colframe=teal!35!black,title={Before you move on}]
What does \fa{عنوان} mean? (\textquotedblleft{}A title; a heading\textquotedblright{}.) Say it once more.
\end{tcolorbox}
\section[\texorpdfstring{\fa{جلد} (jeld)}{جلد (jeld)}]{\fa{جلد} (jeld) — a volume; a (book) cover}

\label{lesson:FA-C528-jeld}



\begin{quote}
Before the new one: say the Persian for a bookshop, then the Persian for a title.
\end{quote}

\subsection*{You'll want to know: \fa{جلد}}
\textbf{\fa{جلد}} (\emph{jeld}) — \textquotedblleft{}a volume; a (book) cover\textquotedblright{}.

\begin{quote}
\fa{این} \fa{رمان} \fa{دو} \fa{جلد} \fa{دارد}.
\end{quote}

\begin{quote}
\emph{in român do jeld dârad.}
\end{quote}

\begin{quote}
This novel is in two volumes.
\end{quote}

Also the counter for books: \emph{se jeld ketâb} 'three books'.

\begin{cousinweb}
\emph{jeld} is a loan, usually traced to Arabic \emph{jild}, which EtymDB glosses 'skin, hide' (EtymDB): the leather a book was bound in, then the book itself. Urdu \emph{jild} is usually traced to the same word. EtymDB has no Persian entry for it.
\end{cousinweb}

\begin{grammarlens}[title={What you've built}]
One more word for telling, explaining and describing at B1.
\end{grammarlens}

\subsection*{Your turn}
\begin{itemize}
\raggedright
  \item \emph{Say it:} \emph{jeld}
  \item \emph{Say it:} \emph{jeld}, once more
  \item \emph{Say it:} a sentence of your own with \emph{jeld}
  \item \emph{Recall:} say the Persian for a bookshop, then the Persian for a title, then say \emph{jeld} again
\end{itemize}

\begin{tcolorbox}[breakable,colback=teal!4,colframe=teal!35!black,title={Before you move on}]
What does \fa{جلد} mean? (\textquotedblleft{}A volume; a (book) cover\textquotedblright{}.) Say it once more.
\end{tcolorbox}
\section[\texorpdfstring{\fa{ناشر} (nâsher)}{ناشر (nâsher)}]{\fa{ناشر} (nâsher) — a publisher}

\label{lesson:FA-C528-nasher}



\begin{quote}
Before the new one: say the Persian for a title, then the Persian for a volume.
\end{quote}

\subsection*{You'll want to know: \fa{ناشر}}
\textbf{\fa{ناشر}} (\emph{nâsher}) — \textquotedblleft{}a publisher\textquotedblright{}.

\begin{quote}
\fa{ناشر} \fa{کتاب} \fa{تازه} \fa{او} \fa{را} \fa{چاپ} \fa{کرد}.
\end{quote}

\begin{quote}
\emph{nâsher ketâb-e tâze-ye u râ châp kard.}
\end{quote}

\begin{quote}
The publisher printed her new book.
\end{quote}

Plural \emph{nâsherân}.

\begin{cousinweb}
\emph{nâsher} is a loan, usually traced to Arabic \emph{nāšir}, 'one who spreads abroad, a publisher'. EtymDB has no entry for it.
\end{cousinweb}

\begin{grammarlens}[title={What you've built}]
One more word for telling, explaining and describing at B1.
\end{grammarlens}

\subsection*{Your turn}
\begin{itemize}
\raggedright
  \item \emph{Say it:} \emph{nâsher}
  \item \emph{Say it:} \emph{nâsher}, once more
  \item \emph{Say it:} a sentence of your own with \emph{nâsher}
  \item \emph{Recall:} say the Persian for a title, then the Persian for a volume, then say \emph{nâsher} again
\end{itemize}

\begin{tcolorbox}[breakable,colback=teal!4,colframe=teal!35!black,title={Before you move on}]
What does \fa{ناشر} mean? (\textquotedblleft{}A publisher\textquotedblright{}.) Say it once more.
\end{tcolorbox}
\section[\texorpdfstring{\fa{مترجم} (motarjem)}{مترجم (motarjem)}]{\fa{مترجم} (motarjem) — a translator}

\label{lesson:FA-C528-motarjem}



\begin{quote}
Before the new one: say the Persian for a volume, then the Persian for a publisher.
\end{quote}

\subsection*{You'll want to know: \fa{مترجم}}
\textbf{\fa{مترجم}} (\emph{motarjem}) — \textquotedblleft{}a translator\textquotedblright{}.

\begin{quote}
\fa{مترجم} \fa{این} \fa{کتاب} \fa{را} \fa{از} \fa{انگلیسی} \fa{ترجمه} \fa{کرد}.
\end{quote}

\begin{quote}
\emph{motarjem in ketâb râ az engelisi tarjome kard.}
\end{quote}

\begin{quote}
The translator translated this book from English.
\end{quote}

Plural \emph{motarjemân}. It shares its root with \emph{tarjome} 'translation'.

\begin{cousinweb}
\emph{motarjem} comes from Arabic \emph{mutarjim}, 'translator' (EtymDB), from the same root as \emph{tarjome}, 'translation'.
\end{cousinweb}

\begin{grammarlens}[title={What you've built}]
One more word for telling, explaining and describing at B1.
\end{grammarlens}

\subsection*{Your turn}
\begin{itemize}
\raggedright
  \item \emph{Say it:} \emph{motarjem}
  \item \emph{Say it:} \emph{motarjem}, once more
  \item \emph{Say it:} a sentence of your own with \emph{motarjem}
  \item \emph{Recall:} say the Persian for a volume, then the Persian for a publisher, then say \emph{motarjem} again
\end{itemize}

\begin{tcolorbox}[breakable,colback=teal!4,colframe=teal!35!black,title={Before you move on}]
What does \fa{مترجم} mean? (\textquotedblleft{}A translator\textquotedblright{}.) Say it once more.
\end{tcolorbox}
